\documentclass[../main.tex]{subfiles}

\begin{document}

\onlyinsubfile{
    \renewcommand{\contentsname}{ОГЛАВЛЕНИЕ}
    \setcounter{tocdepth}{3}
    \tableofcontents
}

\chapter*{\MakeUppercase{Введение}}
\addcontentsline{toc}{chapter}{Введение}

Разработка новых лекарственных препаратов является одной из наиболее сложных и ресурсоемких задач современной науки. Традиционные подходы к дизайну молекул часто связаны с длительным процессом перебора и синтеза кандидатов. В последние годы методы глубокого обучения произвели революцию во многих областях, и хемоинформатика не стала исключением. Применение генеративных моделей, способных изучать закономерности в больших массивах химических данных и создавать новые, химически осмысленные соединения, открывает беспрецедентные возможности для ускорения процесса открытия лекарств \cite{merk2018novo, mouchlis2021advances}. Использование нейронных сетей, в частности рекуррентных (RNN) и архитектур на их основе, позволяет рассматривать молекулярные структуры как своего рода «химический язык», что дает возможность применять для их генерации мощные подходы из области обработки естественного языка \cite{ertl2017silico}. Исследование и совершенствование этих методов является ключевым фактором для создания эффективных инструментов автоматизированного дизайна молекул с заданными биологическими и физико-химическими свойствами.

Несмотря на значительный прогресс, применение генеративных моделей в дизайне лекарств сталкивается с рядом фундаментальных проблем. Во-первых, это проблема адекватного представления молекулярной структуры в формате, понятном для нейронной сети. Линейные нотации, такие как SMILES, удобны, но вносят синтаксические и семантические ограничения, что приводит к генерации невалидных или химически неадекватных структур \cite{wigh2022review}. Во-вторых, существует проблема управляемости генеративного процесса. Базовые модели способны генерировать разнообразные молекулы, но не позволяют целенаправленно создавать соединения с заранее определенным набором свойств. Решение этой проблемы требует разработки более сложных архитектур, таких как вариационные автоэнкодеры (VAE) и модели на основе обучения с подкреплением (RL), каждая из которых имеет свои сильные и слабые стороны.


\begin{itemize}
    \item \textbf{Объект исследования} -- процесс \textit{de novo} дизайна молекулярных структур с помощью методов глубокого обучения.
    \item \textbf{Предмет исследования} -- методы представления (в частности, SMILES) и генерации (на основе RNN, VAE и RL) молекулярных строк для целенаправленного создания химических соединений с желаемыми свойствами.
\end{itemize}

Целью данной курсовой работы является комплексное исследование и сравнительный анализ современных методов представления и генерации молекулярных строк на основе нейронных сетей для задач целенаправленного дизайна молекул.

Для достижения поставленной цели были определены следующие задачи:
\begin{enumerate}
    \item Изучить теоретические основы представления молекулярных структур для задач машинного обучения, уделив особое внимание нотации SMILES, ее преимуществам и недостаткам.
    \item Рассмотреть базовые архитектуры генеративных моделей, такие как рекуррентные нейронные сети и вариационные автоэнкодеры, и их раннее применение для генерации химических структур.
    \item Провести детальный анализ современных подходов к целевой генерации, включая оптимизацию в латентном пространстве (гетероэнкодеры, условные VAE) и методы обучения с подкреплением (RL).
    \item Исследовать проблемы стандартизации оценки генеративных моделей, фундаментальные ограничения строковых представлений и определить перспективные направления развития в данной области.
\end{enumerate}

В ходе работы применялись следующие методы: системный анализ научной и технической литературы по темам хемоинформатики и глубокого обучения; сравнительный анализ архитектур нейронных сетей и подходов к целевой генерации; методы статического анализа кода для изучения практических реализаций моделей; обобщение и систематизация данных для формирования целостной картины текущего состояния и перспектив развития области.

Работа носит обзорно-аналитический характер. Ее практическая значимость заключается в систематизации знаний о современных архитектурах и подходах к генеративному дизайну молекул. Представленный анализ может служить теоретической основой для выбора и разработки эффективных вычислительных инструментов для \textit{de novo} создания новых лекарственных препаратов и химических соединений с заданными свойствами.

\end{document}